% ID: SEL-D0080-OA-E03
% ESTADO: cotejada. Segunda pasada visual y borrador compilado revisados.
% Fuente DOC: Selectividad/sel_2012_matematicas_aplicadas/2 MATEMÁT. APLIC. CC. SS/EXAMEN 5-ANDALUCÍA.doc
% SHA-256 original: 87b729cdd4632386690a9390bdfa17dd7d5c902a2cf7e5801feea45e2b181969
% Página convertida: 1; opción A; ejercicio 3.

Una empresa dispone de tres máquinas A, B y C, que fabrican, respectivamente, el $60\%$, $30\%$ y $10\%$ de los artículos que comercializa.

El $5\%$ de los artículos que fabrica A, el $4\%$ de los de B y el $3\%$ de los de C son defectuosos. Elegido, al azar, un artículo de los que se fabrican en la empresa:
\begin{enumerate}[label=\alph*)]
\item \textbf{(0.5 puntos)} \textquestiondown{}Cuál es la probabilidad de que sea defectuoso y esté fabricado por la máquina C?
\item \textbf{(1.25 puntos)} \textquestiondown{}Cuál es la probabilidad de que no sea defectuoso?
\item \textbf{(0.75 puntos)} Si sabemos que no es defectuoso, \textquestiondown{}cuál es la probabilidad de que proceda de la máquina A?
\end{enumerate}
